\section{总结与延伸}

\subsection{一条可迁移到其他生成系统的主线}

本讲最可迁移的方法不是复制 Movie Gen 的模型名，而是按依赖顺序设计系统：先定义要生成的随机变量，再选择运输目标；先核算 token 长度，再选择 attention 与并行；先定义评估维度，再看样例和 scaling；最后把应用拆成基础生成、编辑、身份和音频等不同条件问题。

\subsection{三个开放研究问题}

第一，能否在不牺牲细粒度运动的前提下进一步压缩时空 latent 或使用稀疏/分层 attention，从而突破长视频二次成本？第二，媒体 reasoning 应产生何种中间状态，怎样构造与人类判断一致的 verifier？第三，原生音视频联合模型能否利用跨模态信息，同时避免对齐数据稀缺和任务互相干扰？

这些问题共同说明，下一代视频模型不会只靠扩大 denoiser。更可能的突破来自 representation、data、attention、parallelism、verification 和 multimodal objective 的协同设计。评价新系统时，最有价值的问题不是“它能生成多炫的视频”，而是“它在哪个 workload、哪组证据和哪条成本边界上改变了 Pareto frontier”。

\subsection{拓展阅读}

\begin{itemize}[itemsep=0.2em,topsep=0.2em]
  \item 官方录像：\url{https://www.youtube.com/watch?v=YGHF8_tf--g}
  \item CS25 V5 课程页：\url{https://web.stanford.edu/class/cs25/past/cs25-v5/}
  \item Movie Gen 一手论文：\url{https://ai.meta.com/static-resource/movie-gen-research-paper/}
  \item 延伸阅读方向：Flow Matching、DiT/AdaLN、Temporal Autoencoder、context parallelism、human evaluation 与 media reasoning verification。
\end{itemize}

\end{document}
